% !TEX program = xelatex
% Compile with: tectonic -X compile moore57_d19_contradiction.tex
\documentclass[11pt,a4paper]{article}

\usepackage{amsmath,amssymb,amsthm,mathtools}
\usepackage[margin=25mm]{geometry}
\usepackage{xcolor}
\usepackage{hyperref}
\usepackage{enumitem}

% --- XeLaTeX / CJK ---
\usepackage{fontspec}
\usepackage{xeCJK}
\setCJKmainfont[BoldFont={Hiragino Mincho ProN W6}]{Hiragino Mincho ProN}
\setCJKsansfont{Hiragino Sans W4}
\setCJKmonofont{Hiragino Sans W4}
\xeCJKsetup{CJKecglue={\hskip 0.15em plus 0.05em minus 0.05em}}

% --- 行間調整 ---
\linespread{1.15}

% --- 定理環境 ---
\theoremstyle{plain}
\newtheorem{theorem}{定理}[section]
\newtheorem{proposition}[theorem]{命題}
\newtheorem{lemma}[theorem]{補題}
\newtheorem{corollary}[theorem]{系}

\theoremstyle{definition}
\newtheorem{definition}[theorem]{定義}
\newtheorem*{fact*}{既知事実}
\newtheorem{fact}[theorem]{既知事実}
\newtheorem{notation}[theorem]{記法}

\theoremstyle{remark}
\newtheorem*{remark}{注意}

% --- 証明環境を「証明」に ---
\renewcommand{\proofname}{\textbf{証明}}

% --- hyperref 設定 ---
\hypersetup{
  colorlinks=true,
  linkcolor=blue!50!black,
  urlcolor=blue!60!black,
  citecolor=blue!50!black,
  pdftitle={Moore graph of degree 57 における D_{19} 作用の高次整合性と矛盾},
  pdfauthor={Yawara Ishida}
}

% --- マクロ ---
\DeclareMathOperator{\Fix}{Fix}
\DeclareMathOperator{\Aut}{Aut}
\DeclareMathOperator{\Sym}{Sym}
\DeclareMathOperator{\Tr}{Tr}

\title{Moore graph of degree $57$ における $D_{19}$ 作用の高次整合性と矛盾}
\author{Yawara Ishida\thanks{\texttt{yawara@aisr.dev}}}
\date{2026年5月13日}

\begin{document}
\maketitle

\begin{abstract}
本稿では、仮想的な Moore graph $\Gamma$ of degree $57$ and diameter $2$ に、位数 $38$ の二面体群
\[
D_{19}=\langle r,t\mid r^{19}=t^2=1,\ trt=r^{-1}\rangle
\]
が自己同型群として作用する、あるいは少なくとも部分群として作用する、という仮定から矛盾を導く。

使う既知事実は次の二つである。
\begin{enumerate}[label=(\arabic*)]
  \item Moore graph of diameter $2$ は、次数が $2,3,7,57$ の場合に限って算術的に可能であり、次数 $57$ の場合は未解決である。
  \item Moore graph of degree $57$ の自己同型群が involution $t$ を含むなら、$\Fix(t)$ は $56$ 頂点からなる star である。
\end{enumerate}

主な新しい部分は、固定点 $u$ のまわりの $57$ 本の枝と $57\cdot 56$ 個の葉を、$D_{19}$-作用と枝間完全マッチングの整合性まで追跡するところにある。一次的な軌道分解
\[
1^1,\quad 19^{55},\quad 38^{58}
\]
だけでは矛盾しない。しかし、$A$-側 fiber 間の「中点反射」条件と、$B,C$-側の $4$-cycle 禁止を合わせると矛盾する。

結論は次である。
\begin{quote}
仮想的 Moore graph $\Gamma$ of degree $57$ and diameter $2$ は、自己同型群に $D_{19}$ を部分群として含まない。特に、$\Aut(\Gamma)\simeq D_{19}$ は不可能である。
\end{quote}

\noindent 注意：本稿の主張は、下記既知事実を前提にした自前の整合性解析である。文献中の既成定理として引用しているものではない。
\end{abstract}

\bigskip
\hrule
\bigskip

\section{既知事実と基本記法}

\begin{fact}[Moore graph of degree $57$]
仮想的 Moore graph $\Gamma$ of degree $57$ and diameter $2$ は、存在すれば強正則グラフ
\[
(v,k,\lambda,\mu)=(3250,57,0,1)
\]
である。したがって、任意の隣接 $2$ 頂点は共通近傍を持たず、任意の非隣接 $2$ 頂点は一意の共通近傍を持つ。

特に、$\Gamma$ は三角形も $4$-cycle も持たない。$4$-cycle があれば、対角の非隣接 $2$ 頂点が $2$ 個の共通近傍を持ち、$\mu=1$ に反するからである。

隣接行列を $A$ とすると、
\[
A^2 = 56I - A + J
\]
であり、スペクトルは
\[
57^1,\qquad 7^{1729},\qquad (-8)^{1520}
\]
である。
\end{fact}

\begin{fact}[Moore graph の可能次数]
Hoffman--Singleton 型の定理により、diameter $2$ の Moore graph の次数は
\[
2,3,7,57
\]
に限られる。したがって次数 $19$ または $38$ の Moore graph は存在しない。
\end{fact}

\begin{fact}[involution の固定点集合]
Moore graph $\Gamma$ of degree $57$ の自己同型群が involution $t$ を含むなら、
\[
\Fix(t)
\]
は $56$ 頂点からなる star である。すなわち、ある中心頂点 $c$ と、その $55$ 個の固定近傍からなる
\[
K_{1,55}
\]
である。
\end{fact}

\begin{notation}
自己同型 $g$ に対して、
\[
a_0(g) = |\Fix(g)|, \qquad
a_1(g) = |\{x \in V(\Gamma) : x \sim g(x)\}|
\]
とおく。
\end{notation}

\bigskip
\hrule
\bigskip

\section{位数 $19$ の回転はただ $1$ 点を固定する}

\begin{lemma}[固定部分グラフの次数は $0,19,38,57$ のいずれかである]\label{lem:degH}
$r$ を位数 $19$ の自己同型とし、
\[
H = \Fix(r)
\]
を $r$ の固定頂点集合で誘導される部分グラフとする。

任意の $x \in H$ について、$r$ は $x$ の $57$ 個の近傍を置換する。固定されない近傍は $19$ 周期に分かれるので、$H$ 内での $x$ の次数 $d_H(x)$ は
\[
57 - d_H(x) \equiv 0 \pmod{19}
\]
を満たす。ゆえに
\[
d_H(x) \in \{0,19,38,57\}.
\]
\end{lemma}

\begin{proof}
$r$ の近傍集合 $N(x)$ 上の作用の軌道長は $1$ または $19$ である。長さ $1$ の軌道が、ちょうど $H$ 内の近傍である。よって $57 - d_H(x)$ は $19$ の倍数である。$0 \le d_H(x) \le 57$ から結論が従う。
\end{proof}

\begin{lemma}[$|H|>1$ なら $H$ は正則 Moore graph である]\label{lem:HMoore}
$|H|>1$ と仮定する。このとき $H$ は正則であり、その次数は
\[
19,38,57
\]
のいずれかである。さらに、次数が $d$ なら $H$ は Moore graph of diameter $2$ and degree $d$ であり、
\[
|H| = d^2 + 1
\]
である。
\end{lemma}

\begin{proof}
まず、$H$ は辺を持つ。もし $x,y \in H$ が非隣接なら、$\mu = 1$ により $x,y$ の一意の共通近傍 $z$ が存在する。$x,y$ は $r$ で固定されるため、一意性から $z$ も $r$ で固定される。ゆえに $z \in H$ であり、$H$ には辺がある。

次に正則性を示す。$x,y \in H$ が非隣接なら、$x,y$ の一意の共通近傍を $z$ とする。任意の
\[
a \in N_H(x) \setminus \{z\}
\]
について、$a$ と $y$ は非隣接である。もし $a \sim y$ なら、$x,y$ は $a$ と $z$ という $2$ つの共通近傍を持つからである。したがって、$a,y$ の一意の共通近傍が存在し、それは $N_H(y) \setminus \{z\}$ に入る。この対応は逆向きにも同様に作れ、単射性も一意性から従う。よって
\[
d_H(x) = d_H(y)
\]
である。

次に $x \sim y$ の場合を考える。補題~\ref{lem:degH} より、$x$ と $y$ の $H$ 内次数は $0$ ではありえず、少なくとも $19$ である。したがって
\[
z \in N_H(x) \setminus \{y\}, \qquad w \in N_H(y) \setminus \{x\}
\]
を取れる。三角形がないので $z \not\sim y$ および $w \not\sim x$ である。また $z \sim w$ なら
\[
x - z - w - y - x
\]
が $4$-cycle になるので不可能である。よって
\[
z \not\sim y, \qquad z \not\sim w, \qquad w \not\sim x.
\]
非隣接頂点同士の次数一致を用いて、
\[
d_H(x) = d_H(w) = d_H(z) = d_H(y)
\]
を得る。したがって $H$ は正則である。

次数を $d$ とする。補題~\ref{lem:degH} と $|H|>1$ から
\[
d \in \{19,38,57\}.
\]
$H$ は三角形を持たず、任意の非隣接 $2$ 頂点は $H$ 内に一意の共通近傍を持つ。したがって $H$ は直径 $2$、次数 $d$ の Moore graph であり、固定頂点 $x$ から数えると
\[
|H| = 1 + d + d(d-1) = d^2 + 1.
\]
\end{proof}

\begin{proposition}[$r$ の固定点はただ $1$ 点である]\label{prop:fix1}
位数 $19$ の自己同型 $r$ について、
\[
|\Fix(r)| = 1.
\]
\end{proposition}

\begin{proof}
頂点数は
\[
3250 \equiv 1 \pmod{19}
\]
なので、$r$ の固定点数は $1$ modulo $19$ である。したがって少なくとも $1$ 点は固定する。

もし $|H|>1$ なら、補題~\ref{lem:HMoore} により $H$ は次数 $19,38,57$ のいずれかの Moore graph である。次数 $19$ または $38$ の Moore graph は存在しない。次数 $57$ の場合、$H$ の任意の頂点の全近傍が固定されるので、$\Gamma$ の連結性から $r$ は全頂点を固定し、$r = 1$ となる。これは $r$ の位数が $19$ であることに反する。

よって $|H| = 1$ である。
\end{proof}

\bigskip
\hrule
\bigskip

\section{固定点 $u$ のまわりの枝と葉}

命題~\ref{prop:fix1} により、$r$ の唯一の固定点を
\[
u
\]
と書く。$\langle r \rangle$ は $D_{19}$ の正規部分群であるから、この唯一の固定点 $u$ は $D_{19}$ 全体で固定される。

\begin{definition}[枝と fiber]
$u$ の近傍集合を $N(u)$ と呼び、その元を「枝」と呼ぶ。

各枝 $b \in N(u)$ に対し、
\[
L_b = \{x \in V(\Gamma) : d(u,x) = 2,\ x \sim b\}
\]
とおく。$L_b$ を枝 $b$ の fiber と呼ぶ。

$\mu = 1$ より、各距離 $2$ 頂点 $x$ は $u$ と一意の共通近傍を持つ。したがって
\[
V(\Gamma) = \{u\} \sqcup N(u) \sqcup \bigsqcup_{b \in N(u)} L_b,
\]
かつ
\[
|N(u)| = 57, \qquad |L_b| = 56.
\]
\end{definition}

\begin{lemma}[異なる fiber 間には完全マッチングが入る]\label{lem:matching}
異なる枝 $b,c \in N(u)$ に対し、$L_b$ と $L_c$ の間の辺は完全マッチングをなす。
\end{lemma}

\begin{proof}
$x \in L_b$ を取る。$x$ は $u$ と非隣接であり、$u,x$ の一意の共通近傍は $b$ である。したがって $x$ は $c$ と隣接しない。よって $x,c$ は非隣接であり、$\mu = 1$ により一意の共通近傍を持つ。その共通近傍は $c$ に隣接し、かつ $u$ とは距離 $2$ にあるので、$L_c$ に属する。

従って、任意の $x \in L_b$ は $L_c$ にちょうど $1$ 個の隣接点を持つ。同様に逆向きも成り立つので、$L_b$ と $L_c$ の間の辺は完全マッチングである。
\end{proof}

\begin{proposition}[枝集合上の $D_{19}$-作用]\label{prop:D19branches}
添字を取り直すと、
\[
N(u) = A \sqcup B \sqcup C
\]
と書け、
\begin{gather*}
A = \{a_i : i \in \mathbb{Z}/19\}, \quad
B = \{b_i : i \in \mathbb{Z}/19\}, \quad
C = \{c_i : i \in \mathbb{Z}/19\}, \\
r(a_i) = a_{i+1}, \qquad r(b_i) = b_{i+1}, \qquad r(c_i) = c_{i+1}, \\
t(a_i) = a_{-i}, \qquad t(b_i) = c_{-i}, \qquad t(c_i) = b_{-i}
\end{gather*}
が成り立つ。
\end{proposition}

\begin{proof}
$r$ は $N(u)$ 上に固定点を持たない。もし枝 $b \in N(u)$ を固定すれば、$b$ は $r$ の固定頂点になり、命題~\ref{prop:fix1} に反するからである。よって $57$ 本の枝は $r$ の $3$ 個の $19$ 周期に分かれる。

反射 $t$ は $u$ を固定する。$\Fix(t)$ は $K_{1,55}$ である。もし $u$ がこの star の中心なら、$N(u)$ のうち $55$ 個の枝が $t$ で固定される。しかし $t$ は各 $19$ 周期上では高々 $1$ 点しか固定できないので、これは不可能である。

したがって $u$ は固定 star の葉であり、固定 star の中心は $N(u)$ 内の枝である。その枝を $a_0$ とする。すると $t$ は $A$ を反転し、
\[
t(a_i) = a_{-i}
\]
となる。

残り $2$ つの $19$ 周期について、$t$ がそれぞれを保存すると、それぞれに $1$ 個ずつ固定枝が生じる。しかし固定 star の中心は $a_0$ であり、他の固定枝は $a_0$ に隣接していなければ固定 star に入れない。枝同士はともに $u$ に隣接するため、三角形禁止により互いに隣接しない。よってこれは不可能である。

したがって残り $2$ つの $19$ 周期は $t$ により交換される。添字を整えることで
\[
t(b_i) = c_{-i}, \qquad t(c_i) = b_{-i}
\]
を得る。
\end{proof}

\bigskip
\hrule
\bigskip

\section{$A$-fiber 間の完全マッチングと中点反射}

以後、
\[
L_i := L_{a_i}
\]
と書く。

\begin{definition}[$A$-fiber の座標化]\label{def:coords}
$L_0$ を固定し、その $56$ 点集合を $P$ と同一視する。$r^i$ により、各 $L_i$ を
\[
L_i = \{(i,p) : p \in P\}
\]
と書く。すなわち
\[
r(i,p) = (i+1, p).
\]

$t$ の $L_0$ 上の作用を
\[
\theta : P \to P
\]
と書く。$t$ の固定 star は中心 $a_0$ を持ち、葉として $u$ と $L_0$ 内の $54$ 点を含む。したがって
\[
P = F \sqcup E, \qquad |F| = 54, \qquad E = \{+,-\},
\]
かつ
\[
\theta|_F = \mathrm{id}, \qquad \theta(+) = -, \qquad \theta(-) = +
\]
である。

この座標化では
\[
t(i,p) = (-i, \theta p).
\]

さらに、
\[
t_h := r^h t r^{-h}
\]
とおくと、
\[
t_h(i,p) = (2h - i, \theta p).
\]

特に $t_h$ は $L_h$ を保存し、$L_h$ 内では $F$ 型の $54$ 点を固定し、$E$ 型の $2$ 点を交換する。
\end{definition}

\begin{definition}[完全マッチングを表す置換 $A_d$]
$d \in \mathbb{Z}/19$, $d \neq 0$ に対し、$L_i$ と $L_{i+d}$ の間の完全マッチングを置換
\[
A_d \in \Sym(P)
\]
で表す。すなわち、
\[
(i,p) \sim (i+d, A_d(p))
\]
と定める。

無向性と反射対称性から、
\[
A_{-d} = A_d^{-1} = \theta A_d \theta
\]
が成り立つ。
\end{definition}

\begin{lemma}[中点反射条件]\label{lem:midpoint}
任意の $h \neq 0$ と $p \in P$ について、
\[
A_{2h}(p) = \theta p \quad \Longleftrightarrow \quad A_h(p) \in E.
\]
\end{lemma}

\begin{proof}
$x = (0,p)$ とおく。中点反射 $t_h$ は
\[
x = (0,p)
\]
を
\[
t_h x = (2h, \theta p)
\]
へ送る。

まず、$x$ と $t_h x$ が隣接する条件は、$L_0$ と $L_{2h}$ のマッチングの定義より、
\[
A_{2h}(p) = \theta p
\]
である。

次に、$x$ の $L_h$ 内の隣接点は
\[
(h, A_h(p))
\]
である。一方、$t_h x$ の $L_h$ 内の隣接点は、上の辺を $t_h$ で写すことで
\[
(h, \theta A_h(p))
\]
である。したがって、$x$ と $t_h x$ が $L_h$ 内に共通近傍を持つことと、
\[
A_h(p) = \theta A_h(p)
\]
は同値である。これは
\[
A_h(p) \in F
\]
と同値である。

もし $A_h(p) \in F$ なら、$x$ と $t_h x$ は共通近傍を持つ。隣接していれば三角形ができるので、隣接しない。

もし $A_h(p) \in E$ なら、$L_h$ 内に固定共通近傍はない。仮に $x$ と $t_h x$ が非隣接なら、$\mu = 1$ により一意の共通近傍 $z$ が存在する。$t_h$ は $x$ と $t_h x$ を交換するので、一意性から $z$ は $t_h$ で固定される。$t_h$ の固定 star の中心は $a_h$ であり、固定頂点は
\[
\{a_h, u\} \cup \{(h,f) : f \in F\}
\]
である。しかし $x, t_h x$ は $u$ にも $a_h$ にも隣接しないので、共通近傍は $L_h$ 内の固定点でなければならない。これは存在しない。よって $x$ と $t_h x$ は隣接する。

以上より、
\[
A_{2h}(p) = \theta p \quad \Longleftrightarrow \quad A_h(p) \in E
\]
である。
\end{proof}

\begin{definition}[例外集合 $S_h$]
$h \neq 0$ に対して
\[
S_h := A_h^{-1}(E)
\]
とおく。明らかに
\[
|S_h| = 2.
\]
補題~\ref{lem:midpoint} は、
\[
S_h = \{p \in P : A_{2h}(p) = \theta p\}
\]
と書き換えられる。また、
\[
S_{-h} = \theta(S_h)
\]
が成り立つ。
\end{definition}

\begin{lemma}[例外集合は $E$ と交わらない]\label{lem:ShE}
任意の $h \neq 0$ について、
\[
S_h \cap E = \varnothing.
\]
\end{lemma}

\begin{proof}
\[
e_h := |S_h \cap E|
\]
とおく。$p \in S_h \cap E$ なら、補題~\ref{lem:midpoint} の形
\[
A_{2h}(p) = \theta p
\]
より、$A_{2h}(p) \in E$ である。したがって
\[
p \in S_{2h}.
\]
よって
\[
S_h \cap E \subseteq S_{2h} \cap E,
\]
したがって
\[
e_h \le e_{2h}.
\]
$\mathbb{Z}_{19}^{\ast}$ 上で $h \mapsto 2h$ は $1$ つの巡回置換をなす。実際、
\[
2^9 \equiv -1 \pmod{19},
\]
なので $2$ の位数は $18$ である。従って、すべての $e_h$ は同じ値 $e$ を持つ。

\medskip
\noindent\textbf{$e = 1$ の排除.}\quad
ある $p \in E$ が $S_h \cap E$ に入るとする。包含関係と濃度一定性から、同じ $p$ が
\[
S_h, S_{2h}, S_{4h}, \ldots
\]
すべてに入る。特に $2^9 h = -h$ なので、
\[
p \in S_{-h}.
\]
一方、
\[
S_{-h} = \theta(S_h)
\]
であり、$S_h \cap E = \{p\}$ なら
\[
S_{-h} \cap E = \{\theta p\}.
\]
$\theta p \neq p$ だから矛盾である。

\medskip
\noindent\textbf{$e = 2$ の排除.}\quad
この場合、すべての $h \neq 0$ について
\[
S_h = E
\]
である。補題~\ref{lem:midpoint} より、任意の $d \neq 0$ について
\[
A_d(+) = -, \qquad A_d(-) = +.
\]
したがって、任意の $j \neq 0$ について、$(0,+)$ と $(j,+)$ は隣接しない。

ところが、任意の $k \neq 0, j$ に対し、$(k,-)$ は $(0,+)$ と $(j,+)$ の共通近傍である。実際、
\[
A_k(+) = -, \qquad A_{k-j}(+) = -
\]
だからである。これは少なくとも $17$ 個の共通近傍を与え、$\mu = 1$ に反する。

したがって $e = 0$ であり、
\[
S_h \cap E = \varnothing
\]
が従う。
\end{proof}

\begin{corollary}[$A$-側 fiber からの回転隣接寄与は常に $38$ である]\label{cor:Aside}
任意の $d \neq 0$ について、$A$-側の葉
\[
\bigsqcup_i L_i
\]
から $a_1(r^d)$ への寄与はちょうど
\[
38
\]
である。
\end{corollary}

\begin{proof}
$x = (i,p) \in L_i$ について、
\[
r^d x = (i+d, p).
\]
したがって $x \sim r^d x$ であることと、
\[
A_d(p) = p
\]
は同値である。

$d = 2h$ と書く。補題~\ref{lem:midpoint} と例外集合 $S_h$ の定義より、
\[
S_h = \{p : A_d(p) = \theta p\}.
\]
補題~\ref{lem:ShE} により $S_h \subset F$ であるから、$p \in S_h$ では $\theta p = p$ であり、従って
\[
A_d(p) = p.
\]
逆に $A_d(p) = p$ なら、$p \notin E$ である。なぜなら $S_d = A_d^{-1}(E)$ は補題~\ref{lem:ShE} により $E$ と交わらず、$p \in E$ なら $A_d(p) = p \in E$ となって矛盾するからである。よって $p \in F$ であり、$\theta p = p$ なので $p \in S_h$ である。

したがって $A_d$ の固定点はちょうど $S_h$ の $2$ 点である。各 fiber $L_i$ で $2$ 点、$i \in \mathbb{Z}/19$ が $19$ 個あるので、寄与は
\[
19 \cdot 2 = 38
\]
である。
\end{proof}

\bigskip
\hrule
\bigskip

\section{跡公式と表現論的制約}

\begin{lemma}[Higman 型跡公式]\label{lem:Higman}
任意の自己同型 $g$ について、$7$-固有空間上の指標値を $\chi_7(g)$ と書くと、
\[
\chi_7(g) = \frac{8 a_0(g) + a_1(g) - 65}{15}
\]
である。
\end{lemma}

\begin{proof}
$7$-固有空間への射影は
\[
E_7 = \frac{A + 8I}{15} - \frac{J}{750}
\]
である。実際、$A$ の固有値 $57,7,-8$ に対して、この作用はそれぞれ $0,1,0$ になる。

$g$ の置換行列を $P_g$ とする。すると
\[
\chi_7(g) = \Tr(E_7 P_g).
\]
ここで
\[
\Tr(P_g) = a_0(g), \qquad
\Tr(A P_g) = a_1(g),
\]
かつ $J P_g = J$ なので
\[
\Tr(J P_g) = \Tr(J) = 3250.
\]
従って
\[
\chi_7(g) = \frac{a_1(g) + 8 a_0(g)}{15} - \frac{3250}{750} = \frac{8 a_0(g) + a_1(g) - 65}{15}.
\]
\end{proof}

\begin{lemma}[反射の $a_0, a_1$ 値]\label{lem:reflection}
反射 $t$ について、
\[
a_0(t) = 56, \qquad a_1(t) = 112.
\]
したがって
\[
\chi_7(t) = 33.
\]
\end{lemma}

\begin{proof}
$a_0(t) = 56$ は既知事実(involution の固定点集合)である。

固定 star の中心を $c$ とする。$c$ は固定葉 $55$ 個と隣接しているので、固定集合の外へ出る辺は
\[
57 - 55 = 2
\]
本である。各固定葉は、固定集合内では中心 $c$ のみに隣接するので、固定集合の外へ
\[
57 - 1 = 56
\]
本の辺を出す。よって固定集合から非固定集合へ出る辺数は
\[
2 + 55 \cdot 56 = 3082
\]
である。

非固定頂点は
\[
3250 - 56 = 3194
\]
個なので、$t$ の $2$ 点軌道は
\[
3194 / 2 = 1597
\]
個である。

$2$ 点軌道 $\{x, tx\}$ が非隣接なら、$\mu = 1$ により一意の共通近傍が存在し、それは $t$ で固定される。逆に、固定頂点から非固定 $2$ 点軌道へ出る辺は、その非隣接 $2$ 点軌道に対して $2$ 本ずつ現れる。したがって、非隣接な $2$ 点軌道の数は
\[
3082 / 2 = 1541
\]
である。

従って、隣接している $2$ 点軌道の数は
\[
1597 - 1541 = 56.
\]
$a_1(t)$ はその両端を数えるので、
\[
a_1(t) = 2 \cdot 56 = 112.
\]

補題~\ref{lem:Higman} より、
\[
\chi_7(t) = \frac{8 \cdot 56 + 112 - 65}{15} = 33.
\]
\end{proof}

\begin{lemma}[頂点置換表現の分解]\label{lem:perm-decomp}
$D_{19}$ の有理既約表現を
\[
\mathbf{1}, \qquad \varepsilon, \qquad \rho
\]
とする。ここで $\mathbf{1}$ は自明表現、$\varepsilon$ は反射に $-1$ を与える $1$ 次元表現、$\rho$ は $18$ 次元の有理既約表現であり、
\[
\rho(1) = 18, \qquad \rho(r^i) = -1\ (i \neq 0), \qquad \rho(t r^i) = 0
\]
である。

このとき、頂点集合上の置換表現 $\pi$ は
\[
\pi = 114\,\mathbf{1} + 58\,\varepsilon + 171\,\rho
\]
と分解する。
\end{lemma}

\begin{proof}
まず軌道構造を確認する。
\begin{itemize}
  \item $u$ は $D_{19}$ 全体で固定されるので、サイズ $1$ の軌道が $1$ 個ある。
  \item 枝集合では、$A$ がサイズ $19$ の軌道、$B \cup C$ がサイズ $38$ の軌道である。
  \item $A$-側の葉では、各 $f \in F$ からサイズ $19$ の軌道が $1$ 個ずつ出るので、サイズ $19$ の葉軌道が $54$ 個ある。また $E = \{+,-\}$ は反射で交換されるため、サイズ $38$ の葉軌道が $1$ 個ある。
  \item $B,C$-側の葉は反射で交換され、反射固定点を持たない。したがってサイズ $38$ の軌道が
  \[
  \frac{2 \cdot 19 \cdot 56}{38} = 56
  \]
  個ある。
\end{itemize}

したがって全体の軌道構造は
\[
1^1, \qquad 19^{55}, \qquad 38^{58}.
\]

サイズ $19$ の軌道表現は
\[
\mathbf{1} + \rho
\]
であり、サイズ $38$ の正則軌道表現は
\[
\mathbf{1} + \varepsilon + 2\rho
\]
である。よって
\[
\pi = \mathbf{1} + 55(\mathbf{1} + \rho) + 58(\mathbf{1} + \varepsilon + 2\rho) = 114\,\mathbf{1} + 58\,\varepsilon + 171\,\rho.
\]
\end{proof}

\begin{proposition}[非自明回転の $a_1$ 値の候補]\label{prop:a1cand}
任意の $d \neq 0$ について、
\[
a_1(r^d) \in \{57, 342, 627, 912\}.
\]
\end{proposition}

\begin{proof}
$7$-固有空間上の有理表現を
\[
\chi_7 = \alpha\,\mathbf{1} + \beta\,\varepsilon + \gamma\,\rho
\]
と書く。

補題~\ref{lem:reflection} より
\[
\alpha - \beta = 33.
\]
また $7$-固有空間の次元は $1729$ なので、
\[
\alpha + \beta + 18\gamma = 1729.
\]
非自明回転 $r^d$ については命題~\ref{prop:fix1} より
\[
a_0(r^d) = 1.
\]
補題~\ref{lem:Higman} から
\[
\chi_7(r^d) = \frac{a_1(r^d) - 57}{15}.
\]
一方、有理既約表現の値から
\[
\chi_7(r^d) = \alpha + \beta - \gamma.
\]
さらに、全頂点表現から定数表現と $7$-固有空間を除いた部分が $(-8)$-固有空間である。補題~\ref{lem:perm-decomp} より、
\[
\chi_{-8} = (114\,\mathbf{1} + 58\,\varepsilon + 171\,\rho) - \mathbf{1} - \chi_7.
\]
したがって各係数は非負整数でなければならない。

上の条件を解くと、ある整数 $s$ について
\[
\alpha = 62 + 9s, \qquad \beta = 29 + 9s, \qquad \gamma = 91 - s,
\]
かつ
\[
\chi_7(r^d) = 19s.
\]
よって
\[
\frac{a_1(r^d) - 57}{15} = 19s,
\]
すなわち
\[
a_1(r^d) = 57 + 285s.
\]

$(-8)$-固有空間側の係数は
\[
(51 - 9s)\,\mathbf{1} + (29 - 9s)\,\varepsilon + (80 + s)\,\rho
\]
である。非負性と $a_1(r^d) \ge 0$ より、
\[
s = 0,1,2,3.
\]
したがって
\[
a_1(r^d) \in \{57, 342, 627, 912\}.
\]
\end{proof}

\begin{corollary}[実際には $a_1(r^d) = 342$ または $912$]\label{cor:a1real}
任意の $d \neq 0$ について、
\[
a_1(r^d) = 342 \quad\text{または}\quad 912.
\]
さらに、$B,C$-側のサイズ $38$ 軌道のうち、$r^d$ に対して寄与する軌道数を $N_d$ とすると、
\[
N_d = 8 \quad\text{または}\quad 23.
\]
特に
\[
N_d \ge 8 \qquad (d \neq 0)
\]
である。
\end{corollary}

\begin{proof}
系~\ref{cor:Aside} より、$A$-側の葉からの寄与は常に $38$ である。枝そのものは $u$ に共通に隣接するため互いに隣接しないので、枝からの寄与はない。また $u$ 自身も寄与しない。

$B,C$-側の葉はサイズ $38$ の $D_{19}$-軌道 $56$ 個に分かれる。各軌道は、ある $d$ に対して寄与するなら $38$ 頂点すべてで寄与し、寄与しないなら $0$ である。

従って
\[
a_1(r^d) = 38 + 38 N_d
\]
と書ける。

命題~\ref{prop:a1cand} の候補
\[
57, 342, 627, 912
\]
のうち、$38$ の倍数であるのは
\[
342, \qquad 912
\]
だけである。従って
\[
a_1(r^d) = 342 \quad\text{または}\quad 912.
\]
このとき
\[
N_d = \frac{a_1(r^d) - 38}{38}
\]
なので、
\[
N_d = 8 \quad\text{または}\quad 23.
\]
\end{proof}

\bigskip
\hrule
\bigskip

\section{$B,C$-側の $4$-cycle 障害}

\begin{definition}[$B,C$-側のサイズ $38$ 軌道]\label{def:BCorbit}
$B,C$-側の葉軌道を $1$ つ固定し、
\[
Q_q = \{b_i^q, c_i^q : i \in \mathbb{Z}/19\}
\]
と書く。ただし
\[
r(b_i^q) = b_{i+1}^q, \qquad r(c_i^q) = c_{i+1}^q,
\]
かつ
\[
t(b_i^q) = c_{-i}^q
\]
となるように記号を選ぶ。

この軌道の $B$-側について、
\[
D_q := \{d \in \mathbb{Z}_{19}^{\ast} : b_i^q \sim b_{i+d}^q\ \text{for all } i\}
\]
と定める。

無向性により、
\[
d \in D_q \quad\Longrightarrow\quad -d \in D_q.
\]
また、$d \in D_q$ であることは、この軌道 $Q_q$ が $a_1(r^d)$ に寄与することと同値である。
\end{definition}

\begin{lemma}[各 $Q_q$ は高々 $1$ つの符号対しか持たない]\label{lem:Dqbound}
任意の $q$ について、
\[
|D_q| \le 2.
\]
\end{lemma}

\begin{proof}
もし
\[
a, b \in D_q, \qquad a \neq \pm b
\]
なら、次の $4$ 頂点を考える:
\[
b_0^q, \qquad b_a^q, \qquad b_{a+b}^q, \qquad b_b^q.
\]
これらは相異なる。実際、$a, b \neq 0$ であり、$a = b$ も $a = -b$ も仮定に反する。

$d = a$ と $d = b$ が $D_q$ に属するので、次の辺が存在する:
\[
b_0^q \sim b_a^q, \qquad b_a^q \sim b_{a+b}^q, \qquad b_{a+b}^q \sim b_b^q, \qquad b_b^q \sim b_0^q.
\]
したがって
\[
b_0^q - b_a^q - b_{a+b}^q - b_b^q - b_0^q
\]
は $4$-cycle である。しかし $\Gamma$ は $4$-cycle を持たない。矛盾。

従って $D_q$ は高々 $1$ つの符号対 $\{d, -d\}$ しか含めず、
\[
|D_q| \le 2
\]
である。
\end{proof}

\bigskip
\hrule
\bigskip

\section{主定理}

\begin{theorem}[$D_{19}$ 部分群は存在しない]\label{thm:main}
仮想的 Moore graph $\Gamma$ of degree $57$ and diameter $2$ について、
\[
D_{19} \le \Aut(\Gamma)
\]
は不可能である。
\end{theorem}

\begin{proof}
系~\ref{cor:a1real} により、任意の非零差
\[
d \in \mathbb{Z}_{19}^{\ast}
\]
について、$B,C$-側のサイズ $38$ 軌道のうち $a_1(r^d)$ に寄与する軌道数 $N_d$ は
\[
N_d \ge 8
\]
を満たす。

一方、定義~\ref{def:BCorbit} より、
\[
N_d = |\{q : d \in D_q\}|.
\]
したがって
\[
\sum_{q} |D_q| = \sum_{d \in \mathbb{Z}_{19}^{\ast}} N_d \ge 18 \cdot 8 = 144.
\]

しかし、$B,C$-側のサイズ $38$ 軌道は $56$ 個であり、補題~\ref{lem:Dqbound} により各軌道について
\[
|D_q| \le 2
\]
である。従って
\[
\sum_q |D_q| \le 56 \cdot 2 = 112.
\]

これは
\[
144 \le 112
\]
という矛盾である。

よって
\[
D_{19} \le \Aut(\Gamma)
\]
は不可能である。
\end{proof}

\begin{corollary}[自己同型群が位数 $38$ の二面体群である場合も不可能]
\[
\Aut(\Gamma) \simeq D_{19}
\]
は不可能である。
\end{corollary}

\begin{proof}
もし $\Aut(\Gamma) \simeq D_{19}$ なら、当然 $D_{19} \le \Aut(\Gamma)$ である。これは定理~\ref{thm:main} に反する。
\end{proof}

\bigskip
\hrule
\bigskip

\section{検証メモ}

\subsection*{7.1\quad 一次軌道分解だけでは矛盾しない}

$D_{19}$ 作用から得られる頂点軌道構造
\[
1^1, \qquad 19^{55}, \qquad 38^{58}
\]
は、頂点数
\[
1 + 55 \cdot 19 + 58 \cdot 38 = 3250
\]
と整合する。したがって矛盾は、単純な軌道長の割り算からは出ない。

\subsection*{7.2\quad 反射 fixed star の中心は $u$ ではない}

$u$ は全 $D_{19}$ で固定されるが、反射の固定 star の中心ではない。もし $u$ が中心なら、反射は $N(u)$ の $55$ 点を固定する必要がある。しかし $N(u)$ は $3$ つの $19$ 周期に分かれ、反射は各周期で高々 $1$ 点しか固定できない。

したがって $u$ は固定 star の葉であり、中心はある枝 $a_i$ である。この点を誤ると、誤った即時矛盾が出る。

\subsection*{7.3\quad 中点反射条件の注意点}

補題~\ref{lem:midpoint} では、$(0,p)$ の $L_h$ 内隣接点は
\[
(h, A_h(p))
\]
であるが、$t_h(0,p)$ の $L_h$ 内隣接点は
\[
(h, \theta A_h(p))
\]
である。したがって、共通近傍が存在する条件は
\[
A_h(p) = \theta A_h(p),
\]
すなわち
\[
A_h(p) \in F
\]
である。この点を明示して検算した。

\subsection*{7.4\quad 最終矛盾の数え上げ}

$A$-側の寄与は各非零 $d$ について常に $38$ である。従って、表現論的候補
\[
57, 342, 627, 912
\]
のうち、全体が $38$ の倍数になる候補だけが残る。これが
\[
342, \qquad 912
\]
であり、$B,C$-側には各 $d$ につき少なくとも $8$ 個のサイズ $38$ 軌道が必要になる。

一方、各 $B,C$ 軌道は $4$-cycle 禁止により高々 $1$ つの符号対 $\{d,-d\}$ にしか寄与できない。したがって全体の容量は
\[
56 \cdot 2 = 112
\]
であるが、必要量は
\[
18 \cdot 8 = 144
\]
である。この不等式衝突が最終矛盾である。

\bigskip
\hrule
\bigskip

\section*{参考文献}

\begin{itemize}\setlength{\itemsep}{4pt}
  \item M.~Miller and J.~Širáň, \textit{Moore graphs and beyond: A survey of the degree/diameter problem}, Electronic Journal of Combinatorics, Dynamic Survey DS14.\\
  \url{https://www.combinatorics.org/files/Surveys/ds14/ds14v2-2013.pdf}\\
  Diameter $2$ の Moore graph の可能次数、および degree $57$ の未解決性について参照。
  \item A.~A.~Makhnev, \textit{Partial geometries and their extensions}, Russian Mathematical Surveys 53:6.\\
  \url{https://www.mathnet.ru/php/getFT.phtml?jrnid=rm\&paperid=203\&what=fullteng}\\
  Proposition 1.3 に、degree $57$ の Aschbacher graph に involution がある場合、固定点集合が $56$ 頂点の star であることが述べられている。
  \item A.~E.~Brouwer and W.~H.~Haemers, \textit{Spectra of Graphs}.\\
  \url{https://homepages.cwi.nl/~aeb/math/ipm/ipm.pdf}\\
  Higman method、association scheme の自己同型に対する跡公式、および Moore graph of degree $57$ に関する標準的制約について参照。
\end{itemize}

\end{document}
